% Target: rmp-workbench-iii-diagram-two
% Formula: Closing the virtual boundary of Y followed by A gives Y.
% Ink: Kernel flat frame; the Y ring stands over the A atom and four glyphless
%   corners carry the rectangular right-hand virtual return.
% Boundary: The virtual word is traced; the two physical indices of A are open.
% Source: author source ImagesReview Section III/ImagesReview Section III.tex
%   lines 164-169; archival output Dia2.pdf
% Capabilities: kernel, periodic-closure, boundary-insertion
\[
  \tenkzkernel
  \begin{tenkz}[rows={wire,wire,wire,wire}, cols=2, bonds=none]
    \tn[at=(2,1), skin=ring, name=Y,
      ports={90:virtual, 270:virtual}]{Y}
    \tn[at=(3,1), skin=mpo, name=A,
      ports={0:physical, 90:virtual, 180:physical, 270:virtual}]{A}
    \tn[at=(1,1), skin=none, name=NW, ports={0:virtual, 270:virtual}]{}
    \tn[at=(1,2), skin=none, name=NE, ports={180:virtual, 270:virtual}]{}
    \tn[at=(4,1), skin=none, name=SW, ports={0:virtual, 90:virtual}]{}
    \tn[at=(4,2), skin=none, name=SE, ports={90:virtual, 180:virtual}]{}
    \tnwire{NW.270}{Y.90} \tnwire{Y.270}{A.90} \tnwire{A.270}{SW.90}
    \tnwire{SW.0}{SE.180} \tnwire{SE.90}{NE.270} \tnwire{NE.180}{NW.0}
  \end{tenkz} \;:\;Y
\]
